\section{Experimental Methodology}
\label{sec:methodology}

\subsection{\textsc{skewbench}}

We implement the study as an open harness rather than a set of scripts, because
the artifact is as much the contribution as the numbers. \textsc{skewbench}
comprises a workload generator, six join strategies, an event-log metrics
parser, an experiment driver, and the cost model of
Section~\ref{sec:costmodel}. Every result in this paper is produced by a single
command and every figure and table is generated from the results file with no
manual transcription.

\subsection{Workload generation}

The generator produces two Parquet tables with an aerospace engine-telemetry
schema, following the NASA C-MAPSS turbofan degradation
convention~\cite{saxena2008cmapss} --- unit identifier, operational cycle,
three operational settings, and $N_s$ sensor channels --- so that the synthetic
generator can be exchanged for the real dataset without touching the benchmark.

\texttt{fact\_telemetry} is the skewed side: one row per sensor snapshot, with
\texttt{engine\_id} drawn from a truncated Zipf law over $n$ engines. This
mirrors a real fleet, in which a small number of problem units generate a
disproportionate share of recorded events. \texttt{dim\_maintenance} is the
probe side: one row per maintenance work order. Its keys are uniform except for
the hot key, which carries $\mu$ times the mean per-key row count.

The choice to parameterise two-sidedness by $\mu$ rather than by drawing the
probe side from the same Zipf law is the single most consequential design
decision in the harness, and it is not cosmetic. Under a doubly-Zipf design the
hot key alone yields $H \times P$ output rows. With $H \approx 4.4\times10^{5}$
and $P \approx 4.4\times10^{3}$ that is roughly $2\times10^{9}$ rows from a
single key --- a thousandfold amplification of a two-million-row fact table,
which we observed as a run that made no progress and abandoned. Such a workload
does not measure skew mitigation; it measures the cost of materialising a
near-cartesian product, and every strategy looks equally bad because all of
them are dominated by output volume. Bounding $\mu$ keeps the
output tractable and promotes $P$ from an artefact of the generator to a swept
independent variable, which is what makes the $1/\sqrt{P}$ prediction testable
at all.

Skewed variants of TPC-H exist~\cite{microsoft_skewgen, tpch} and the harness
accepts them, but they do not provide a controllable $\mu$, which is the
condition our decision rule turns on. Generation is deterministic given a seed,
and the manifest written alongside the data records $H$ and $P$ exactly, so no
later stage must re-derive them.

\subsection{Strategies evaluated}

\begin{enumerate}
\item \textbf{baseline} --- sort-merge join, AQE disabled, no mitigation.
\item \textbf{aqe} --- AQE enabled with skew-join optimisation.
\item \textbf{salt\_uniform}$(k)$ --- whole probe side replicated $k$ times,
AQE disabled.
\item \textbf{salt\_selective}$(k)$ --- hot-key rows only replicated, AQE
disabled.
\item \textbf{aqe\_salt}$(k)$ --- AQE enabled \emph{and} selective salting.
\item \textbf{broadcast} --- broadcast hash join, where the probe side fits.
\end{enumerate}

Automatic broadcast is disabled throughout
(\texttt{autoBroadcastJoinThreshold} $=-1$) except in the broadcast arm, so
that a strategy under test is never silently replaced by a different plan. Each
arm's executed physical plan is captured into the results file, so that the
claim ``this was a sort-merge join'' is auditable rather than asserted.

\subsection{Metrics}

We separate metrics by what they license.

\textbf{Deterministic.} Shuffle bytes read and written, records shuffled, and
bytes spilled to memory and disk are properties of the physical plan and the
partitioning. They are identical on a laptop and on a hundred-node cluster
given the same plan and the same data, and the quantitative claims about
replication cost rest on them.

\textbf{Structural.} The \emph{skew ratio} --- maximum over median task
duration within the join stage --- and the task-time coefficient of variation
are within-run ratios. They cancel most machine-specific scaling while still
capturing the straggler that skew produces, and they are how we measure whether
a mitigation actually worked as opposed to merely running faster.

\textbf{Machine-dependent.} Wall-clock and executor run time. Reported, used to
calibrate $\gamma$, and never load-bearing for a transferable claim.

Metrics are recovered by parsing the Spark event log --- a JSON-lines file
written by Spark itself --- rather than through a listener callback. The log
survives the driver exiting and can be re-analysed later without re-running
anything, which is what makes a result auditable months after the fact. Each
measurement is tagged with a unique Spark job group, giving an exact rather
than timestamp-inferred mapping from a results row to the tasks that produced
it.

\subsection{Execution protocol}

Each cell runs one warmup execution, discarded, followed by five timed
repetitions. Warmup matters because the first execution of a plan pays JIT
compilation, class loading and page-cache misses, and including it inflates the
mean unevenly across arms. We report the \emph{median} and interquartile range
rather than mean and standard deviation, because run times on a shared machine
have a long right tail.

Materialisation is forced with Spark's \texttt{noop} sink, so no rows are
written to disk or shipped to the driver and the measurement is of the join and
its aggregation alone. The terminal operation is a grouped aggregation touching
columns from both sides, not a \texttt{count()}: counting invites the optimizer
to prune the payload and, on some plans, to skip the probe entirely, which
would measure column pruning rather than the strategy under test.

\subsection{Platform}

Measurements were taken on a single node with \CORES\ cores and \DRIVERMEM\ of
driver memory, Apache Spark \SPARKVER\ on OpenJDK \JAVAVER, with dedicated
local compute --- not shared serverless. This choice is deliberate: Nurdin et
al.~\cite{nurdin2026lakehouseperf} measure roughly twofold run-to-run variance
for identical lakehouse queries on shared cloud compute, which would make
timing claims meaningless. We are trading absolute scale for measurement
control, and Section~\ref{sec:threats} states what that costs us.

The full configuration grid, the raw event logs, and the exact software
versions are included in the artifact.
